\scanpage{207}
\mspage{63}
au foncteur $\alpha_i$, comme

\[
 \operatorname{Hom}\bigl(\alpha_{i!}(\xi_{I_{/i}}),F\bigr)\longrightarrow
 \operatorname{Hom}\bigl(\alpha_{i!}(\mathinner{.\,.\,.\,.}),F\bigr).
\]

Dire que c'est un iso pour $\forall\ F$, signifie donc \uline{exactement} que la transformée de la flèche $(\unsure{xxx|xx})$ par $\alpha_{\unsure{i!|!}}$ est un iso. Pour en conclure que $(xxx)$ elle-$=$ est iso, la condition naturelle à imposer \deleted{\unsure{ici}} est donc que $\alpha_{i!}\colon\mathbb{D}(I_{/i})\to\mathbb{D}(I)$ soit conservatif. C'est le cas $3^\circ$. Cela achève la démonstration de \unsure{4|6}.4.

4.5. Montrons que la condition $3^\circ)$ \inserted{sur $I$} n'est pas tjrs satisfaite, dans le cas de $\mathbb{D}=(\mathrm{HOT}_{W_0})^{\mathrm{opp}}\approx (\mathbb{D}_{\mathrm{Ens}})^{\mathrm{opp}}$, bien que dans ce cas, il soit pourtant \deleted{\unsure{si}} vrai que pour $\forall\ I'$ sur $I$ satisfaisant la condition $a'')$ \inserted{(autoduale)}, la condition $b')$ \inserted{(elle aussi autoduale)} est aussi satisfaite (\inserted{par} cas $4^\circ$ ci-dessus). Je dis donc qu'il existe $I$ et $i$ de $I$, tels que

\[
 \beta_{i*}\colon({}_{i\backslash}I)^\wedge\longrightarrow I^\wedge
\]

ne soit pas conservatif. (\uline{NB} Par contre $\alpha_{i!}\colon(I_{/i})^\wedge\mathbin{\unsure{-|\to}}I^\wedge$ est tjrs conservatif \deleted{\unsure{cf}}, ce qui \deleted{\illegible{1}} redonne \deleted{\unsure{une des raisons}} \inserted{le cas $4^\circ$} comme cas particulier de $3^\circ$)). Comme $\beta_i\colon{}_{i\backslash}I\longrightarrow I$ est $W_0$-propre (et $=$ $W_\infty$-propre !), $\beta_{i*}$ \inserted{\unsure{$F$}} se calcule fibre par fibre'', donc on a
% archive sheet number: 207
